\section{函数 III \& 自然数}

\subsection{逆函数}

\begin{definition}
    设 $f: X \to Y$ 是函数。
    \begin{itemize}
        \item 若存在 $g: Y \to X$ 使得 $g \circ f = I_X$，那么称 $f$ \textbf{左可逆}，称 $g$ 为 $f$ 的一个\textbf{左逆函数}。
        \item 若存在 $g: Y \to X$ 使得 $f \circ g = I_Y$，那么称 $f$ \textbf{右可逆}，称 $g$ 为 $f$ 的一个\textbf{右逆函数}。
        \item 若存在 $g: Y \to X$ 使得 $g \circ f = I_X \land f \circ g = I_Y$，那么称 $f$ \textbf{可逆}，称 $g$ 为 $f$ 的一个\textbf{逆函数}。
    \end{itemize}
\end{definition}

函数可逆性可以用下面的定理判定：

\begin{theorem}
    设 $X, Y$ 为集合且 $X \neq \varnothing$，函数 $f: X \to Y$，那么下面条件等价：
    \begin{enumerate}
        \item $f$ 为单射；
        \item $f$ 左可逆；
        \item $f$ 可左消去，即对任意 $Z$ 以及 $g: Z \to X, h: Z \to X$，有 $f \circ g = f \circ h \implies g = h$。
    \end{enumerate}
\end{theorem}

\begin{theorem}
    设 $X, Y$ 为集合且，函数 $f: X \to Y$，那么下面条件等价：
    \begin{enumerate}
        \item $f$ 为满射；
        \item $f$ 右可逆；
        \item $f$ 可右消去，即对任意 $Z$ 以及 $g: Y \to Z, h: Y \to Z$，有 $g \circ f = h \circ f \implies g = h$。
    \end{enumerate}
\end{theorem}

\begin{theorem}
    设 $X, Y$ 为集合且，函数 $f: X \to Y$，那么下面条件等价：
    \begin{enumerate}
        \item $f$ 为双射；
        \item $f$ 左可逆、右可逆、可逆；
        \item $f$ 的逆关系 $f^{-1}$ 也是 $f$ 的逆函数。
    \end{enumerate}
\end{theorem}

\subsection{集合的特征函数}

\begin{definition}[特征函数]
    设全集 $U$，集合 $A \subseteq U$ 是集合，那么 $A$ 的\textbf{特征函数}定义为：
    $$
    \chi_{A}(x) = \begin{cases}
        1,\ x \in A \\
        0,\ x \notin A \\
    \end{cases}
    $$
\end{definition}

特征函数满足下面的性质：

\begin{itemize}
    \item $0 \le \chi_A \le 1$；
    \item $\chi_A = \mathbf{0} \iff A = \varnothing$；
    \item $\chi_A = \mathbf{1} \iff A = U$；
    \item $\chi_A \le \chi_B \iff A \subseteq B$。
\end{itemize}

\subsection{自然数与数学归纳法}

首先定义集合的后继：

\begin{definition}[后继]
    若 $A$ 为集合，定义 $A \cup \{A\}$ 为 $A$ 的\textbf{后继}。记做 $A^{+}$。
\end{definition}

那么，空集 $\varnothing$ 的后继是 $\{\varnothing\}$，它的后继是 $\{\varnothing, \{\varnothing\}\}$，它的后继是 $\{\varnothing, \{\varnothing\}, \{\varnothing, \{\varnothing\}\}\}$ ……

显然，每个集合都有唯一的后继，且满足 $A \in A^{+},\ A \subseteq A^{+}$。

给出了集合的后继定义，我们可以开始构造自然数。这里给出 Von Neumann 自然数构造方案：
$$
\begin{gathered}
    0 = \varnothing,\ 1 = 0^{+},\ 2 = 1^{+},\ 3 = 2^{+}, \dots, n + 1 = n^{+}, \dots \\
    \NN = \{0, 1, 2, \dots, n, \dots\}
\end{gathered}
$$

然后，我们可以定义自然数的比较：

\begin{definition}[自然数的小于]
    若 $m, n \in \NN$ 使得 $m \in n$，则称 $m$ 小于 $n$，记作 $m < n$。
\end{definition}

那么小于关系是 $\NN$ 上的反自反、对称、传递关系。

接着我们定义自然数的加法运算和乘法运算：

\begin{definition}[自然数的加法和乘法]
    对任意 $n, m \in \NN$，令：
    \begin{itemize}
        \item $m + 0 = m,\ m \cdot 0 = 0$。
        \item $m + n^{+} = (m + n)^{+},\ m \cdot n^{+} = m \cdot n + m$。
    \end{itemize}
\end{definition}

由此，我们得到\textbf{自然数系统} $\langle \NN, +, \cdot \rangle$。

\begin{theorem}[Peano 公理]
    Von Neumann 自然数系统满足 Peano 公理：
    $$
    \begin{aligned}
        & P_1: 0 \in \NN \\
        & P_2: n \in \NN \implies n^{+} \in \NN \\
        & P_3: n \in \NN \implies n^{+} \neq 0 \\
        & P_4: n, m \in \NN \land n^{+} = m^{+} \in \NN \implies n = m \\
        & P_5: 0 \in S \land (n \in S \to n^{+} \in S) \implies S = \NN
    \end{aligned}
    $$
\end{theorem}